\subsection{Videre arbeid våren 2027}

\subsubsection{Trehjulsbil med ett forhjul}
\label{sec:videre-trehjul}
En trehjulsbil står på tre punkter, og hjullasten avhenger da ikke av hvor rett rammen er (seksjon~\ref{sec:disk-kjoretest}). Den trenger også bare ett lager foran, og får ett hjul mindre å akselerere og lavere masse. Ulempene er at bilen velter lettere, at en liten skjevhet i gaffelen til forhjulet gir sideavvik, og at tyngdepunktet må ligge nær bakakselen for at drivhjulene ikke skal slure. Forslaget kan testes med en utskiftbar front, slik at samme bil kjøres med to forhjul og med ett, med minst ti kjøringer per variant.

Et steg videre er en tohjuling med svinghjulet i en forgaffel som kan dreie, der gyroeffekten kan styre gaffelen inn mot helningen, slik en sykkel holder seg oppreist uten fører. Om det kan bli stabilt i denne størrelsen, må modelleres før noe bygges.

\subsubsection{Dynamisk analyse i MAST2007}
Den samme bilen brukes videre i MAST2007 våren 2027, og følgende er utsatt til da:
\begin{itemize}
    \item Utledning av bevegelseslikningene for svinghjul, drivverk og bil, inkludert sluringen ved nedsetting, og validering mot kjøretestene.
    \item Tap ved turtall over \SI{3100}{o/min}, målt med svinghjulet drevet av en drill, og en sensor på drivakselen som kan måle sluringen direkte.
    \item Optimering av utvekslingen med målte tap i hele drivverket.
\end{itemize}
